\documentclass{article}

\begin{document}
    \begin{center}
        \Large \textbf{Huntsman Challenge 0: Einrichten der Ansible Verwaltungsumgebung} \\ [6pt]
        \normalsize
        Autor: Maximilian Jakob Maag B.Sc. \\
        Version: 1.0
    \end{center}

    \section{Intro}
    Die Challenge richtet sich an die Huntsman Fraktion. Aufgabe ist das Einrichten einer zentralen Verwaltungsumgebung für Ansible.
    Sie werden die zentrale Kontrollinstanz sein, die alle von den Rabbits deployten Server verwaltet.
    \\ \\
    Das Ziel dieser Challenge ist es, eine robuste und sichere Ansible-Umgebung aufzubauen, von der aus Sie alle Server Ihres Rabbit-Konkurrenten managen können.

    \section{Aufgaben}

    \subsection{Aufgabe 1}
    Erstellen Sie ein SSH Schlüsselpaar auf Ihrem lokalen Notebook.

    \subsection{Aufgabe 2}
    Deployen Sie mittels Terraform eine VM, die als Ansible Control Node fungiert.
    Ihre VM muss folgende Eigenschaften besitzen:
    \begin{itemize}
        \item Ubuntu 22.04
        \item 2 vCPUs
        \item 4 GB RAM
        \item 20 GB Festplatte
        \item Port 22 offen
        \item Ihr SSH Public Key im authorized\_keys File
    \end{itemize}

    \subsection{Aufgabe 3}
    Installieren Sie auf der Control Node VM Ansible und alle benötigten Abhängigkeiten.
    \\ \\
    \textbf{Hinweis:} Sie können diese Installation auch lokal auf Ihrem Notebook durchführen, falls Sie das bevorzugen.

    \subsection{Aufgabe 4}
    Erstellen Sie auf Ihrem lokalen Notebook oder der Control Node eine Basis-Verzeichnisstruktur für Ansible Projekte:
    \begin{itemize}
        \item \texttt{inventory/} - für Hostgruppen und Variablen
        \item \texttt{playbooks/} - für Ihre Playbooks
        \item \texttt{roles/} - für wiederverwendbare Rollen
        \item \texttt{group\_vars/} - für Gruppenvariablen
        \item \texttt{host\_vars/} - für Host-spezifische Variablen
    \end{itemize}

    \subsection{Aufgabe 5}
    Erstellen Sie ein erstes Testinventar mit mindestens 3 Hosts.
    Diese Hosts können vorerst nur auf dem Papier existieren und müssen keine echten Maschinen sein.
    Organisieren Sie diese Hosts in logischen Gruppen (z.B. webserver, databases, monitoring).

    \subsection{Aufgabe 6}
    Generieren Sie auf der Ansible Control Node ein zusätzliches SSH Schlüsselpaar.
    Der Public Key dieses Schlüsselpaares wird später an Ihren Rabbit-Konkurrenten übergeben, damit dieser ihn in die authorized\_keys seiner VMs eintragen kann.
    \\ \\
    \textbf{Wichtig:} Speichern Sie den Private Key sicher. Sie werden diesen Key benötigen, um auf alle vom Rabbit deployten Server zuzugreifen.
    Es ist unter keinen Umständen erlaubt, den Private Key mit anderen zu teilen oder öffentlich zugänglich zu machen.
    Er darf die Maschine auf dem er generiert wurde nicht verlassen.

\end{document}
